\chapter{Cryptographic Forgiveness}
\label{short:cryptographic-forgiveness}

Could parties show one another that they have deleted a file or a record?
The wish is to make an ending verifiable, so that an agreement to let
something go need not rest entirely on a promise.

But where has the information traveled? A copy, a cache, a screenshot, or
a compromised device may outlast the declared deletion.

Destroying encryption keys offers a narrower direction to investigate. Its
meaning depends on what was stored, where copies exist, and who retained
access. The aspiration remains valuable precisely because its limits must
be faced: what would the parties actually know after the ceremony ends?

\shortlimit{Deleting a key does not establish that every copy of the information is gone.}
\companionref{C}{appendix.C}{An open deletion problem and a narrower key-destruction research direction.}
